\section{Extensions to modulus constraints}

Recall that in the body of the paper, cardinality constraints, either within $\GPtwo$ or the global cardinality constraints added to $\Ctwo$, did not constrain the modulus of a count. So, for example, in $\GPtwo$ we could not state that the number of neighbors of an element with a given property is even. And similarly, we could not add to $\Ctwo$ a  constraint that the total number of elements in a model satisfying a formula is an even number. We stated in the body that we can add such modulus constraints and obtain the same results.

In this appendix we will first indicate how to extend the satisfiability result for $\GPtwo$ to the case when existential quantifiers are allowed in the Presburger quantifiers, i.e., when the formula $P(x)$ is of the form:
\begin{equation*}
P_i(x)\ :=\ 
\exists x_1 \cdots \exists x_n\ \sum_{t \in \intsinterval{m}} \lambda_{t} \cdot \presby{R_t(x, y) \land x \neq y}
\ + \ \sum_{t\in \intsinterval{n}} \lambda_t'\cdot x_t
\ \circledast_i\ \delta.
\end{equation*}
Note that the existential variables captures a modulus equation, as $a \equiv b \bmod p$ can be rewritten as:
$$
\exists x_1\exists x_2\
a+x_1p = b+x_2p
$$

We claim that Theorem~\ref{thm:gptwoexp} still holds under such a generalization.
The proof is basically the same.
The only difference is that in the definition of characteristic system
$\cC^\varphi_\pi(\bfx)$
defined in Definition~\ref{def:characteristic-system} the variables $\bfx$ include all the existential variables.
Similarly for the system $\cQ^\varphi(\bfx)$
defined in Definition~\ref{def:linear-system-Q}.
Then, we observe that
Lemmas~\ref{lem:equiv1},~\ref{lem:equiv2}, and~\ref{lem:equiv3} still hold for such systems
and the satisfiability problem of $\GPtwo$ stays in $\exptime$.

For the extension of $\Ctwo$ with global cardinality and modulus constraints, the situation is more obvious. We have already shown how to get a (modulus-constraint free) semilinear representation for the type spectrum vectors. We can compose this with modulus constraints in addition to linear arithmetic constraints and solve using standard techniques for Presburger arithmetic. 
Rewriting modulus constraint into a single linear equation does not change the number of constraints in each disjunct in the formula $\Psi^{\varphi}(\bfx)$.
We emphasize again that the argument here is completely generic, and the same comment holds to show that for \emph{any} decidability extension of Presburger arithmetic (e.g. B\"uchi Arithemetic or Semenov Arithmetic \cite{semenovpresexp}), $\Ctwo$ with global cardinality constraints of that form is decidable.
